%% LyX 2.1.4 created this file.  For more info, see http://www.lyx.org/.
%% Do not edit unless you really know what you are doing.
\documentclass[a4paper,oneside,brazil,11pt,a4paper,openright,titlepage,usenames,dvipsnames]{book}
\usepackage[utf8]{inputenc}
\usepackage[T1]{fontenc}
\usepackage{lmodern}
\setcounter{secnumdepth}{3}
\setcounter{tocdepth}{3}
\usepackage{array}
\usepackage{verbatim}
\usepackage{calc}
\usepackage{textcomp}
\usepackage{amssymb}
\usepackage{graphicx}

\makeatletter

%%%%%%%%%%%%%%%%%%%%%%%%%%%%%% LyX specific LaTeX commands.
\pdfpageheight\paperheight
\pdfpagewidth\paperwidth

%% Because html converters don't know tabularnewline
\providecommand{\tabularnewline}{\\}

%%%%%%%%%%%%%%%%%%%%%%%%%%%%%% User specified LaTeX commands.
% Classe alternativa, apropriada para impressão frente-verso. Inclui páginas em branco
% de forma que capítulos sempre tenham início na página à direita:
% \documentclass[11pt,a4paper,openright,titlepage]{book}

% Pacotes
\usepackage[brazilian]{babel}
\usepackage{epsfig}
\usepackage{subfigure}
\usepackage{amsfonts}
\usepackage{amsmath}
\usepackage[thmmarks,amsmath]{ntheorem}%\usepackage{amsthm}
\usepackage{boxedminipage}
\usepackage{geometry}
\usepackage{theorem}
\usepackage{fancybox}
\usepackage{fancyhdr}
\usepackage{ifthen}
\usepackage{url}
\usepackage{afterpage}
\usepackage{color}
\usepackage{colortbl}
\usepackage{rotating}
\usepackage{makeidx}
\usepackage{indentfirst}
% Pacotes para adição de figuras do inkscape
\usepackage{graphicx}
\usepackage{import}

% Escolher um dos seguintes formatos:
\usepackage{ft2unb} % segue padrão de fontes do Latex

\makeindex

\makeatother

\usepackage{babel}
\begin{document}
\setcounter{secnumdepth}{3}
\setcounter{tocdepth}{2}
\pagestyle{empty}

\grau{Engenheiro de Controle e Automação}

\tipodemonografia{TRABALHO DE GRADUAÇÃO}

\begin{comment}
Título
\end{comment}


\titulolinhai{TÍTULO DO TRABALHO DIVIDIDO EM}

\titulolinhaii{MAIS DE UMA LINHA PARA TÍTULOS}

\titulolinhaiii{REALMENTE LONGOS COMO ESTE}

\titulolinhaiv{}

\begin{comment}
Autores. Basta retirar o texto totalmente caso não haja um determinado
autor.
\end{comment}


\autori{João da Silva}

\autorii{}

\autoriii{}

\begin{comment}
Membros da banca. Basta retirar o texto totalmente caso não haja um
determinado membro da banca.
\end{comment}


\membrodabancai{Prof. Fulano de Sousa, ENE/UnB}

\membrodabancaifuncao{Orientador}

\membrodabancaii{Prof. Ciclano dos Santos, ENE/UnB}

\membrodabancaiifuncao{Examinador interno}

\membrodabancaiii{Prof. Beltrano da Silva, ENE/UnB}

\membrodabancaiiifuncao{Examinador interno}

\membrodabancaiv{}

\membrodabancaivfuncao{}

\membrodabancav{}

\membrodabancavfuncao{}

\begin{comment}
Data de defesa: mês e ano
\end{comment}


\mes{dezembro}
\ano{2015}

\begin{comment}
Comandos para criar a capa e a página de assinaturas
\end{comment}


\capaprincipal
\capaassinaturas

\begin{comment}
Ficha Catalográfica
\end{comment}


\noindent \textbf{FICHA CATALOGRÁFICA}

\noindent %
\fbox{\begin{minipage}[t]{1\columnwidth}%
JOÃO, DA SILVA

Título do trabalho dividido em mais de uma linha para títulos realmente
longos como este,

\medskip{}


{[}Distrito Federal{]} 2015.

\medskip{}


x, 101p., 297 mm (FT/UnB, Engenheiro, Controle e Automação, 2015).
Trabalho de Graduação \textendash{} Universidade de Brasília.Faculdade
de Tecnologia.

\medskip{}


1. Bla\hfill{}2.Ble\hfill{}

3. Bli

\medskip{}


I. Mecatrônica/FT/UnB\hfill{}II. Título (Série)\hfill{}

%
\end{minipage}}

\noindent \medskip{}


\noindent \textbf{REFERÊNCIA BIBLIOGRÁFICA}

SILVA, JOÃO DA, (2015). Título do trabalho dividido em mais de uma
linha para títulos realmente longos como este. Trabalho de Graduação
em Engenharia de Controle e Automação, Publicação FT.TG-$n^{\circ}022$,
Faculdade de Tecnologia, Universidade de Brasília, Brasília, DF, 101p.

\noindent \bigskip{}


\noindent \textbf{CESSÃO DE DIREITOS}

\noindent AUTOR: João da Silva

TÍTULO DO TRABALHO DE GRADUAÇÃO: Título do trabalho dividido em mais
de uma linha para títulos realmente longos como este.

\noindent \medskip{}


\noindent GRAU: Engenheiro\hfill{}ANO: 2015\hfill{}

\noindent \medskip{}


É concedida à Universidade de Brasília permissão para reproduzir cópias
deste Trabalho de Graduação e para emprestar ou vender tais cópias
somente para propósitos acadêmicos e científicos. O autor reserva
outros direitos de publicação e nenhuma parte desse Trabalho de Graduação
pode ser reproduzida sem autorização por escrito do autor.

\noindent \bigskip{}


\noindent \rule[0.5ex]{1\columnwidth}{1pt}

\noindent João da Silva

\noindent Rua dos Bobos, nº 0, Bairro Feliz.

\noindent 71000-000 Brasília \textendash{} DF \textendash{} Brasil.


\begin{comment}
Dedicatória
\end{comment}


\frontmatter

\begin{comment}
Texto de dedicatória do primeiro autor.
\end{comment}


\dedicatoriaautori{um beijo pra minha mãe, pro meu pai, e pra você}

\begin{comment}
Texto de dedicatória do segundo autor. Caso não tenha um segundo autor,
este texto não será mostrado
\end{comment}


\dedicatoriaautorii{Dedicatória do autor 2}

\begin{comment}
Texto de dedicatória do terceiro autor. Caso não tenha um terceiro
autor, este texto não será mostrado
\end{comment}


\dedicatoriaautoriii{Dedicatória do autor 3}

\begin{comment}
Comando para criar a página de dedicatória
\end{comment}


\dedicatoria

\begin{comment}
Agradecimentos
\end{comment}


\begin{comment}
Texto de agradecimentos do primeiro autor.
\end{comment}


\agradecimentosautori{Agradecimentos!}

\begin{comment}
Texto de agradecimentos do segundo autor. Caso não tenha um segundo
autor, este texto não será mostrado.
\end{comment}


\agradecimentosautorii{A inclusão desta seção de agradecimentos é
opcional e fica à critério do(s) autor(es), que caso deseje(em) inclui-la
deverá(ão) utilizar este espaço, seguindo esta formatação.}

\begin{comment}
Texto de agradecimentos do terceiro autor. Caso não tenha um terceiro
autor, este texto não será mostrado.
\end{comment}


\agradecimentosautoriii{A inclusão desta seção de agradecimentos
é opcional e fica à critério do(s) autor(es), que caso deseje(em)
inclui-la deverá(ão) utilizar este espaço, seguindo esta formatação.}

\begin{comment}
Comando para criar a página de agradecimentos
\end{comment}


\agradecimentos

\resumo{resumo}{Resumo!

\medskip{}


Palavras Chave: bla, ble, bli

}\vspace*{2cm}


\resumo{Abstract}{Abstract, in English ofc!

\medskip{}


Keywords: bla, ble, bli

}

\begin{comment}
Listas de conteúdo, figuras e tabelas
\end{comment}


\sumario
\listadefiguras
\listadetabelas

\begin{comment}
Lista de Símbolos
\end{comment}


%TCIDATA{LaTeXparent=0,0,these.tex}


%\chapter*{\setfontarial\mdseries LISTA DE SÍMBOLOS} % se usar ft1unb.sty, descomente esta linha



\chapter*{LISTA DE SÍMBOLOS}

% se usar ft2unb.sty, descomente esta linha



\subsection*{Símbolos Latinos}

\begin{tabular}{p{0.1\textwidth}p{0.63\textwidth}>{\PreserveBacklash\raggedleft}p{0.15\textwidth}}
$v$  & Velocidade linear  & {[}m/s{]}\tabularnewline
\end{tabular}


\subsection*{Símbolos Gregos}

\begin{tabular}{p{0.1\textwidth}p{0.63\textwidth}>{\PreserveBacklash\raggedleft}p{0.15\textwidth}}
$\omega$ & Velocidade angular & {[}rad/s{]}\tabularnewline
\end{tabular}


\subsection*{Grupos Adimensionais}

\begin{tabular}{p{0.1\textwidth}p{0.8\textwidth}}
i, k & Contador\tabularnewline
\end{tabular}


\subsection*{Subscritos}

\begin{tabular}{p{0.1\textwidth}p{0.8\textwidth}}
$ref$  & referência \tabularnewline
$fer$  & ferramenta \tabularnewline
$sis$  & sistema \tabularnewline
$des$  & desejado\tabularnewline
\end{tabular}


\subsection*{Sobrescritos}

\begin{tabular}{p{0.1\textwidth}p{0.8\textwidth}}
$\cdot$  & Variação temporal \tabularnewline
$-$  & Valor médio \tabularnewline
\end{tabular}


\subsection*{Siglas}

\begin{tabular}{p{0.1\textwidth}p{0.8\textwidth}}
PCI  & \textit{Peripheral Component Interconnect}\tabularnewline
CPU & Unidade Central de Processamento - \textit{Central Processing Unit} \tabularnewline
AO & Saída Analógica - \textit{Analog Out}\tabularnewline
DO & Saída Digital - \textit{Digital Out}\tabularnewline
CS & Seletor de \textit{Chip - Chip Select}\tabularnewline
SC & Sem Conexão\tabularnewline
P.I. & Placa de Interface\tabularnewline
ICW & \textit{Initialization Command Words}\tabularnewline
OCW & \textit{Operational Control Word}\tabularnewline
\end{tabular}


\begin{comment}
Corpo Principal
\end{comment}


\mainmatter
\setcounter{page}{1}
\pagenumbering{arabic}
\pagestyle{plain}

\begin{comment}
Introdução
\end{comment}

\chapter{Introdução}

\label{CapIntro}

% Resumo opcional. Comentar se não usar.
\resumodocapitulo{Resumo opcional}


\section{Contextualização}

Contextualizar.

Conforme \cite{article:dummy}, vide a Tabela \ref{tab:Descrever-tabela}.
Assim sendo, observe a Figura \ref{fig:Descrever-figura.}.
\begin{table}[h]
\begin{centering}
\begin{tabular}{|c|c|c|c|c|}
\hline
 &  &  &  & \tabularnewline
\hline
\hline
 &  &  &  & \tabularnewline
\hline
 &  &  &  & \tabularnewline
\hline
 &  &  &  & \tabularnewline
\hline
 &  &  &  & \tabularnewline
\hline
\end{tabular}
\par\end{centering}

\label{tab:Descrever-tabela}Descrever tabela.


\end{table}


\begin{figure}[h]
\begin{centering}
\includegraphics[width=0.4\columnwidth]{figs/capa_fundo}
\par\end{centering}

\label{fig:Descrever-figura.}Descrever figura.


\end{figure}



\section{Definição do problema}

Definir problema.


\section{Objetivos do projeto}

Objetivos.


\section{Resultados obtidos}

Resultados.


\section{Apresentação do manuscrito}

Apresentar.


\begin{comment}
Fundamentos
\end{comment}

\chapter{Fundamentos\label{chap:FundamentacaoMatematica}}

% Resumo opcional. Comentar se não usar.
\resumodocapitulo{Resumo opcional.}


\section{Introdução}

Introduzir.


\begin{comment}
Conclusões
\end{comment}
\chapter{Conclusões}\label{CapConclusoes}

Neste trabalho foi feito o estudo da análise estabilidade de sistemas não-lineares por meio do Teorema de estabilidade de Lyapunov via LMIs. Inicialmente considerou-se o sistema com dinâmica não-linear, a partir do qual se obteve o retrato de fase do sistema para a região no espaço de estados. A partir do retrato de fase foi possível verificar qualitativamente o comportamento dos pontos de equilíbrio do sistema, podendo-se classificá-los, dentre outros, em estável, assintoticamente estável e instável. O ponto de equilíbrio é estável quando as trajetórias de resposta para diferentes pontos iniciais sempre se aproximam para uma região próxima ao ponto de equilíbrio. Quando as trajetórias tendem ao ponto de equilíbrio à medida que o tempo tende ao infinito, este ponto é dito como assintoticamente estável. Caso as trajetórias se afastem, o ponto é dito como instável. Este princípio foi utilizado no decorrer do trabalho para se classificar analisar a estabilidade de pontos de equilíbrio.

Em seguida, foi obtido o modelo fuzzy Takagi-Sugeno dos sistemas não-lineares em estudo e verificou-se que esta modelagem representa o modelo linearizado de forma exata. Obtida esta modelagem, que consiste no conjunto de sistemas linearizados associados na forma de vértices, iniciou-se tópico principal deste projeto, que consistiu na análise de estabilidade e na estimativa da região de atração para pontos de equilíbrio na origem. Utilizou-se o artificio de mudança de variável para deslocar o ponto de equilíbrio do sistema para a origem, sem que houvesse perda de generalidade, de forma a se abranger um número maior de sistemas.

No capítulo sobre análise de estabilidade utilizou-se o Teorema de Estabilidade de Lyapunov, LMIs e o Teorema de Finslar para se propor um novo Teorema de análise de estabilidade de sistemas fuzzy T-S modelados segundo a técnica de não-linearidade por setor local. O objetivo com a proposta deste Teorema foi obter condições menos conservadoras, tais que a análise de estabilidade garantisse resultados mais abrangentes. Fez-se uso de outros métodos já existentes na literatura para serem comparados com o método aqui proposto.

As comparações do método proposto neste trabalho com os métodos encontrados na literatura permitiram ver que o nosso método gera melhores resultados para um valor fixo de $\lambda = 20$, visto que foi o único que garantiu a estabilidade em toda a região $C$ para a qual o sistema é definido. Fixando-se a região de análise em todo o domínio $C$, variou-se o valor de $\lambda$ até se obter os valores tais que as LMIs necessárias para a estabilidade não gerassem resultado. Assim, verificou-se que o nosso método foi o melhor quando analisado o limitante superior, uma vez que fora o único para o qual se verificou que não há limitante superior de $\lambda$ para o qual a origem deixe de ser estável para a região $C$. Já para o limitante inferior, o nosso método gerou o pior resultado.

Mantendo-se $\lambda = 20$, foi feito uma estimativa da região de atração para os métodos utilizados até agora, utilizando-se o mesmo critério de maximização para todos. Verificou-se que o nosso método gera o melhor resultado para esta estimativa. Foi sugerido um novo teorema para a maximização da estimativa da região de atração, o qual não foi bem sucedido.

Como trabalhos futuros sugere-se:
\begin{enumerate}
\item Comparar com métodos mais recentes na literatura;
\item Alterar os valores dos pesos $\omega_1$ e $\omega_2$ do Teorema \ref{th:reg_atrac_1} para se obter uma melhor estimativa de região de atração por meio deste;
\item Aumentar os graus dos simplexes para os vértices P e X do resultado principal, em busca de melhores resultados.
\end{enumerate}

\begin{comment}
Bibliografia
\end{comment}


\renewcommand{\bibname}{REFERÊNCIAS BIBLIOGRÁFICAS}
\addcontentsline{toc}{chapter}{REFERÊNCIAS BIBLIOGRÁFICAS}

\bibliographystyle{abnt-num}
\bibliography{bibliography}


\begin{comment}
Anexos
\end{comment}


\anexos
\makeatletter
% não retirar estes comandos
\renewcommand{\@makechapterhead}[1]{%
  {\parindent \z@ \raggedleft \setfontarial\bfseries
\LARGE \thechapter. \space\space
\uppercase{#1}\par
\vskip 40\p@
}
}
\makeatother

\begin{comment}
Anexo I: Descrição do CD
\end{comment}


%TCIDATA{LaTeXparent=0,0,relatorio.tex}

\chapter{Descri��o do conte�do do CD}\label{AnCD}


\refstepcounter{noAnexo}

\begin{comment}
Anexo II: Programas Utilizados
\end{comment}


%TCIDATA{LaTeXparent=0,0,relatorio.tex}

\chapter{\textit{Plot} Retrato de Fase\label{AnEsquematicos}

x1 = x1\_min:x1\_step:x1\_max;

x2 = x2\_min:x2\_step:x2\_max;

[X1,X2] = meshgrid(x1,x2);

F1= f1(X1, X2);

F2= f2(X2, X2);

hold off

figure;clf

streamslice(X1,X2,F1,F2,2); \% https://www.mathworks.com/help/matlab/ref/streamslice.html

xlabel('x\_1');ylabel('x\_2');

line([0],[0],'marker','o','linestyle','none','markerfacecolor','r') \%plots origin as equilibrium point

hold off;

\chapter{Plot Condi��o de Inclus�o (Estimativa de regi�o de atra��o) para P constante\label{AnEsquematicos

function level\_curve(P, gamma, color)

\%Plots the level curve V(x) = x'*P*x = gamma, given P and gamma.

for i = 1:length(P)


if length(P) > 1



Pi = P{i{1;
else
Pi = P{i;
end
%a = Pi(1,1);
b = Pi(1,2); % = Pi(2,1);
c = Pi(2,2);
dP = det(Pi);

%[V,D] = eig(Pi);
%v1=V(:,1);
%v2=V(:,2);
%figure;plot(v1(1),v1(2),'r*',v2(1),v2(2),'b*');
%hold on

%Calculo
x1max=sqrt(c*gamma/dP);
steps = (2*x1max)/1e4;
x1=-x1max:steps:x1max;
x2a=($-b*x1+sqrt(c*gamma-x1.^2*dP))/c$;
x2b=($-b*x1-sqrt(c*gamma-x1.^2*dP))/c$;

%Plot
hold on
plot(x1,x2a,color)
plot(x1,x2b,color)
hold off
grid
volume = det(inv(Pi)*gamma); 
end
end


\refstepcounter{noAnexo}

\begin{comment}
Acrescente mais anexos conforme julgar necessário.
\end{comment}

\end{document}
